% Korean edition of sections/s3_design.tex.
\seclabel
\claim{\secnum{3}{sec:design}\quad 래퍼는 로그인 세션 하나를 유지하고, 전략과 시장 데이터는 호출자에게 남긴다}

\lead{래퍼에는 목표 5개와 비목표 5개가 있으며, 가장 큰 제외 항목은 시장 데이터다(\figref{goals}).} 대회 서버는 포트폴리오가 보유한 종목의 가격만 알려주고(포지션 레코드마다 \code{close}), 동아리의 시장 데이터는 한국투자증권 API에서 온다. 이 API를 래퍼에 넣으면 대회 클라이언트가 증권사의 인증과 요청 한도에 묶이므로, 가격은 호출자가 넘기고 래퍼는 대회 사이트 하나의 클라이언트로 남는다. 그렇다면 \code{plan()}은 보유하지 않은 종목의 주문 크기를 어떻게 정하는가? 호출자가 \code{prices=}로 그 가격을 넘기며, 가격이 없는 주문은 경고와 함께 남겨 둔다(\ref{sec:oneshare}절).

% Korean edition of figs/goals.tex.
\fig[wide]{goals}{목표 5개와 비목표 5개가 래퍼의 범위를 정한다}[\code{pytimefolio}의 범위: 목표마다 달성 여부를 보여주는 기준]{%
{\tabfont\begin{tabularx}{\linewidth}{>{\RaggedRight\arraybackslash}p{7mm} >{\RaggedRight\arraybackslash}p{70mm} >{\RaggedRight\arraybackslash}X}
\thd{ID} & \thd{목표} & \thd{성공 기준}\\
\hl G1 & 사람 없이 대회 기간 내내 로그인 상태를 유지한다 & 시작부터 11월 30일까지 수동 로그인 없음. 모든 갱신과 로그인을 기록\\\rowsep
G2 & 전략에 포트폴리오의 현재 상태를 준다 & (a) 병합: 델타마다 병합 상태가 \code{t}가 같은 스냅샷과 같다. (b) 신선도: 읽기, 주문, \code{wait()}이 26\,s(폴링 2회와 20\,s 응답 마감)보다 오래된 상태를 쓰지 않는다(대신 \code{StaleFeed}). 3\,s 폴링은 목표이지 보장이 아니다\\\rowsep
G3 & 주최측 한도와 대회 가이드라인 안에서 주문을 넣고, 정지하고, 추적한다 & 하루 최대 80건(새 주문과 정지한 주문 합계). 쓰기 사이 최소 1\,s. 서버가 받았을 수 있는 주문은 다시 보내지 않음. 넣은 주문은 15\,s 안에 피드나 REST \code{Orders}에서 확인. \code{plan()}은 계획 시점에 합계, 유니버스, 종목, 섹터 규칙을 어기는 목표 비중을 거부(실행 중 비중은 한도 안에 붙잡아 두지 않음)\\\rowsep
G4 & 폴링이 주문을 늦추지 않는다 & 피드 폴링과 주문이 이벤트 루프 하나를 공유하며, 어떤 주문도 피드 응답을 기다리지 않음(\ref{sec:live}장)\\\rowsep
G5 & 모호하면 추측하지 않고 멈춘다 & 포트폴리오 규칙에 맞는 항목이 없거나 둘 이상일 때, 넣은 주문에 새 레코드가 둘 이상 맞을 때, \code{currBizDay}가 KST 날짜와 다를 때, 넣은 주문을 15\,s 안에 피드에서도 REST \code{Orders}에서도 찾지 못할 때(재전송 안 함) \code{AmbiguousState}. 400이면 \code{RequestError}\\
\specialrule{0.8pt}{1pt}{1pt}
\thd{ID} & \thd{비목표} & \thd{이유}\\
N1 & 시장 데이터(가격, 호가창) & 가격은 호출자가 준다. 다른 API를 넣으면 로그인과 한도가 하나씩 더 생긴다\\\rowsep
N2 & 전략과 포트폴리오 구성 & 래퍼는 목표 비중을 주문으로 바꾼다. 목표를 고르는 것은 호출자 몫이다\\\rowsep
N3 & 브라우저 자동화 & 모든 동작이 페이지 렌더링을 기다리고 레이아웃이 바뀌면 깨진다. 페이지 자체의 HTTP 요청은 16\textasciitilde 39\ms{}가 걸린다(\ref{sec:transport}절)\\\rowsep
N4 & 여러 계정 또는 다른 참가자의 포트폴리오 & 프로세스당 계정 하나. 다른 참가자의 데이터는 우리가 조회할 것이 아니다\\\rowsep
N5 & 요청 한도 탐색 & 주최측은 트래픽 폭주를 이유로 IP 주소를 차단할 수 있다\\
\end{tabularx}}}{G = 목표, N = 비목표. 음영: G1은 11월 30일까지 실행해야만 보일 수 있다. \ref{sec:live}장의 실서버 테스트는 오전 한나절이었다. 80건 한도는 주최측 기준인 하루 100건보다 20건의 여유를 둔다.}{DART Quant Team}


\lead{서버나 주최측이 한도를 정한 곳에서는 래퍼가 그 한도 안쪽으로 여유를 두고 자기 매개변수를 정한다. 주최측 100건 대비 하루 80건, 900\,s 토큰 대비 720\,s 갱신, 15\textasciitilde 30\,s 유휴 종료 대비 \code{keepalive_expiry} 10\,s다.} 20건의 여유는 래퍼가 마지막으로 건수를 동기화한 뒤 브라우저에서 손으로 넣은 주문(건수는 시작할 때와 \code{rebalance()} 때만 서버 수치로 올린다)과, 취소가 100건에 포함되는지라는 미해결 질문(Q2)을 덮는다. 주최측은 권장 요청 간 지연의 수치를 주지 않았다. 1\,s 간격은 쓰기를 초당 1건으로 제한한다. 부록~\ref{app:decisions}는 매개변수마다 그 근거가 된 측정이나 규칙, 그리고 매개변수를 바꾸게 할 조건을 나열한다.

\lead{거래일은 5단계로 진행되며, 로그인과 준비 작업은 모두 영업일 전환(06:02\textasciitilde 07:07)과 08:50 사이에 끝난다(\figref{schedule}).} 장중(09:00\textasciitilde 15:20)에는 래퍼가 실시간 피드를 폴링하고 호출자의 주문을 보내기만 한다. 갱신 타이머는 없다. 서버 시계로 토큰 발급 후 720\,s가 지난 뒤 첫 REST 요청이 토큰을 갱신하고, 실시간 피드는 연결하거나 다시 연결할 때만 토큰을 요청한다. 래퍼는 로그인과 갱신 때만 \code{currBizDay}를 읽으므로 전환 이후에 시작한다.

% Korean edition of figs/schedule.tex.
\fig[wide]{schedule}{래퍼는 08:50 전에 준비를 마치고, 장중에는 폴링과 매매만 한다}[개장 때 한 번 리밸런싱하는 전략의 거래일]{%
{\tabfont\begin{tabularx}{\linewidth}{>{\RaggedRight\arraybackslash}p{20mm} >{\RaggedRight\arraybackslash}p{26mm} >{\RaggedRight\arraybackslash}X}
\thd{시각 (KST)} & \thd{단계} & \thd{래퍼 동작}\\
07:10 & 시작 & 저장된 쿠키 묶음으로 토큰을 갱신하고, 실패하면 로그인. 포트폴리오 선택. \code{currBizDay}가 KST 날짜와 같아질 때까지 주문 거부\\\rowsep
07:15 & 불러오기 & 실시간 피드를 열고 스냅샷을 받음. 종목별 섹터를 위해 \code{Universe}를 불러옴. 서버에서 그날 주문 건수를 셈\\\rowsep
08:00\textasciitilde 08:50 & 계획(호출자) & 호출자가 자체 가격으로 목표 비중을 계산. \code{plan()}이 유니버스, 합계, 종목·섹터 한도, 1주 비중과 대조해 주문으로 바꿈. \code{guidelines()}가 회전율과 소형주 합계를 보고\\\rowsep
08:50 & 제출 & \code{rebalance()}가 서버에서 그날 주문을 다시 세고 하루 한도 80건을 확인한 뒤, 매도부터 최소 1\,s 간격으로 보냄. 각 주문은 피드의 \code{ords}에 나타나면 확인됨\\\rowsep
09:00\textasciitilde 15:20 & 장중 & 3\,s마다 피드를 폴링. 720\,s 뒤 첫 REST 호출이 토큰을 갱신하고, 피드는 재연결할 때만 토큰을 요청. 호출자의 요청에 따라 주문 정지\\
\end{tabularx}}}{시각은 호출자가 정한다. 3\,s 폴링을 빼면 래퍼는 호출될 때만 동작한다. 08:30에 넣은 주문은 08:33에 \code{Working}이 됐고, 장 시작인 09:00에 체결되기 시작해 09:08:07에 \code{Done}이 됐다(\ref{sec:orders}절).}{DART Quant Team}


\lead{페이지가 보내는 것과 같은 HTTP 요청을 직접 보내는 방식을 택했다. 페이지 레이아웃에도 브라우저에도 의존하지 않고, 재사용 연결에서 Python 요청이 31\textasciitilde 39\ms{} 안에 응답받기 때문이다(\ref{sec:transport}절, \figref{alternatives}).} GitHub에서 찾은 다른 대회용 도구들은 Playwright로 웹 페이지를 조작한다. 요청을 역공학하지 않아도 되지만, 모든 동작이 페이지 렌더링을 기다리고, 레이아웃이 바뀌면 스크립트가 깨지며, 전략을 실행하는 컴퓨터에 브라우저를 설치해야 한다. Playwright의 시간은 재지 않았다. 이 판단은 지연 시간 비교가 아니라 레이아웃 의존성과 브라우저 설치에 근거한다.

% Korean edition of figs/alternatives.tex.
\fig[wide]{alternatives}{\code{asyncio} 스레드 하나에서 HTTP 요청을 직접 보내며, 확인되지 않은 주문은 다시 보내지 않는다}[설계 결정마다 검토한 대안]{%
{\tabfont\begin{tabularx}{\linewidth}{>{\RaggedRight\arraybackslash}p{27mm} >{\RaggedRight\arraybackslash}p{52mm} >{\RaggedRight\arraybackslash}X}
\thd{결정} & \thd{선택} & \thd{대안과 탈락 이유}\\
서버 접근 방식 & 페이지가 하는 것과 같은 HTTP 요청: 재사용 요청당 curl 16.0\ms{}, Python 31\textasciitilde 39\ms & 브라우저 자동화(Playwright): 모든 동작이 페이지 렌더링을 기다리고, 레이아웃이 바뀌면 깨지며, 실행 호스트에 브라우저가 필요\\\rowsep
동시성 & 스레드 하나의 \code{asyncio} & 동기 요청: 피드 요청이 응답을 기다리는 동안 주문도 기다리게 된다. 스레드: 피드 스레드가 아무 두 문장 사이에서 상태를 바꿀 수 있어 상태를 읽을 때마다 잠금이 필요하다. \code{asyncio}는 \code{await}에서만 태스크를 바꾼다. 두 태스크 모두 대부분 네트워크를 기다리므로 GIL 비용이 없다\\\rowsep
피드 전송 방식 & WebSocket만 & 서버 협상 응답이 제시하는 Server-Sent Events와 롱 폴링: 롱 폴링은 메시지마다 HTTP 요청이 1개 더 든다. 둘 다 구현하지 않음\\\rowsep
WebSocket의 토큰 & \code{Authorization} 헤더 & 브라우저처럼 URL에 넣기: 토큰이 접근 로그에 남는다\\\rowsep
\code{AddOrder} 응답 유실 & 다시 보내지 않음: 15\,s 동안 피드에서, 다음으로 REST \code{Orders}에서 찾고, 둘 다 없으면 \code{AmbiguousState}. 서버에 닿지 않은 연결 실패 뒤에만 한 번 재전송 & 즉시 재전송: 응답에 주문 ID가 없으므로 주문이 중복될 위험\\
\end{tabularx}}}{‘선택’은 \code{pytimefolio} 0.1.0이 구현한 것이다. 요청 시간은 \ref{sec:transport}절에서 가져왔으며, 브라우저 자동화의 시간은 재지 않았다. 실시간 피드의 협상 응답(\ref{sec:feed}절)은 WebSocket, Server-Sent Events, 롱 폴링을 나열한다.}{DART Quant Team}


\lead{모든 작업은 스레드 하나의 asyncio 이벤트 루프에서 코루틴으로 돌며, 서로 끼어들면 안 되는 단계만 잠금 3개로 보호한다(\figref{asyncio}).} 코루틴은 \code{await}에서만 스레드를 넘기므로, \code{await}가 없는 델타 병합은 다른 코루틴이 반쯤 쓴 상태를 볼 일이 없고, 피드 태스크와 호출자의 코루틴은 잠금 없이 상태를 공유한다. 작업은 대부분 소켓을 기다리는 시간이다. 요청 하나에 31\textasciitilde 39\ms{}가 걸리고, curl의 16.0\ms{} 중 네트워크 왕복만 11.1\ms{}다(\ref{sec:transport}절). 요청당 스레드 자체의 작업 시간은 재지 않았지만, 모든 테스트에서 스레드 하나로 피드와 주문을 따라갔다. 잠금 3개는 중간에 \code{await}가 있는 단계를 덮는다.
\begin{itemize}
\item 갱신은 한 번에 하나만 실행된다. 토큰이 720\,s보다 오래됐음을 안 코루틴은 토큰 잠금을 잡고 그 안에서 나이를 다시 확인한다. 잠금을 기다리던 두 번째 코루틴은 새 토큰을 발견하고 갱신 없이 돌아간다. 오프라인 테스트에서 동시 \code{token()} 호출 5건은 \code{SilentRefresh}를 1번 보냈다.
\item 쓰기는 관문을 거친다. 관문 잠금은 한 번에 쓰기 하나만, 직전 쓰기보다 최소 1\,s 뒤에, 그날 건수에 이번 요청을 더해도 80 이하일 때만 통과시키며, 어떤 요청에서든 403이나 429를 받은 뒤에는 아예 통과시키지 않는다.
\item 주문을 넣을 때는 쓰기와 새 주문 검색 내내 잠금을 잡는다. 주문을 넣는 코루틴은 최대 15\,s 동안 0.5\,s마다 피드 상태에서 주문을 찾고, 피드 자체는 3\,s마다 새로 고쳐진다. \code{place()}가 이미 돌려준 ID는 검색에서 빼므로, 첫 주문이 REST로만 확인된 경우에도 같은 주문 2건이 같은 주문 ID에 묶이지 않는다. 잠금은 이 프로세스에만 적용되므로, 래퍼는 실행 중에 계정당 주문 제출자가 하나라고 가정한다(\ref{sec:orders}절).
\end{itemize}

% Korean edition of figs/asyncio.tex.
% Gantt block: \tfgb{y}{x0}{x1}{fill}{text colour}{label}
\providecommand{\tfgb}[6]{\fill[#4] (#2,#1-2) rectangle (#3,#1+2);
  \node[font=\fontsize{7}{8.4}\selectfont,text=#5,inner sep=0pt] at ({(#2+#3)/2},#1) {#6};}
\fig[wide]{asyncio}{스레드 하나가 코루틴을 번갈아 실행하며, 토큰 잠금은 그중 하나만 갱신하게 한다}[\code{rebalance()}가 주문 2건을 넣고 \code{guidelines()}가 함께 실행되는 동안 이벤트 루프의 간트 차트]{%
\begin{tikzpicture}[x=1mm,y=1mm,
  f/.style={font=\fontsize{7}{8.4}\selectfont},
  lane/.style={f,anchor=east,align=right,text=navy},
  note/.style={f,text=greyD,inner sep=0.5pt},
  ptitle/.style={font=\fontsize{7.5}{9}\selectfont\bfseries,text=navy,anchor=west},
  ev/.style={-{Stealth[length=1.4mm]},line width=0.5pt,navyM,dashed}]
% ====================================================== panel (a): 0 to 145 ms, 1 mm per ms
\node[ptitle] at (-30,8) {(a) 처음 145\,ms: 한 번에 코루틴 하나만 실행된다. G는 다시 갱신하지 않고 토큰 잠금을 기다린다};
\node[lane] at (-2,0)   {\textbf{이벤트 루프}\\{\color{greyD}스레드 하나}};
\node[lane] at (-2,-9)  {\textbf{피드 태스크}\\{\color{greyD}\code{hub.stream()}}};
\node[lane] at (-2,-21) {\textbf{코루틴 R}\\{\color{greyD}\code{rebalance()}}};
\node[lane] at (-2,-32) {\textbf{코루틴 G}\\{\color{greyD}\code{guidelines()}}};
\node[lane] at (-2,-42) {\textbf{토큰 잠금}\\{\color{greyD}\code{Auth._lock}}};
\node[lane] at (-2,-51) {\textbf{주문 잠금}\\{\color{greyD}\code{Orders._place_lock}}};
\begin{scope}
% event loop strip
\draw[greyM,line width=0.4pt] (0,-2) rectangle (145,2);
\tfgb{0}{0}{3}{navy}{white}{R}     \tfgb{0}{3}{6}{navy}{white}{G}
\tfgb{0}{35}{38}{navy}{white}{R}   \tfgb{0}{38}{41}{navy}{white}{G}
\tfgb{0}{70}{73}{navy}{white}{R}   \tfgb{0}{73}{76}{navy}{white}{G}
\tfgb{0}{105}{108}{navy}{white}{R} \tfgb{0}{140}{143}{navy}{white}{R}
\foreach \a/\b in {6/35,41/70,76/105,108/140} {\node[note] at ({(\a+\b)/2},0) {유휴};}
% feed task
\tfgb{-9}{0}{145}{navyL}{navy}{피드 요청 사이 \code{sleep(3)}. 다음 요청은 패널 (b)에}
% coroutine R
\tfgb{-21}{0}{3}{navy}{white}{}
\tfgb{-21}{3}{35}{navyM}{white}{\code{SilentRefresh}}
\tfgb{-21}{35}{38}{navy}{white}{}
\tfgb{-21}{38}{70}{navyM}{white}{\code{Orders} (다시 세기)}
\tfgb{-21}{70}{73}{navy}{white}{}
\tfgb{-21}{73}{105}{navyM}{white}{\code{ValidateProduct}}
\tfgb{-21}{105}{108}{navy}{white}{}
\tfgb{-21}{108}{140}{navyM}{white}{\code{AddOrder}}
\tfgb{-21}{140}{143}{navy}{white}{}
\fill[navyL] (143,-23) rectangle (145,-19);
\draw[decorate,decoration={brace,amplitude=1.2mm},navy,line width=0.5pt] (0,-18.3) -- (36.5,-18.3);
\node[f,text=navy,anchor=south west] at (0,-17.3) {토큰이 720\,s보다 오래됨: 잠금, 재확인, 갱신};
\draw[decorate,decoration={brace,amplitude=1.2mm},navy,line width=0.5pt] (105,-18.3) -- (108,-18.3);
\node[f,text=navy,anchor=south east] at (145,-17.3) {쓰기 관문: 직전 쓰기 후 최소 1\,s, 오늘 80건 미만};
% coroutine G
\tfgb{-32}{3}{6}{navy}{white}{}
\tfgb{-32}{6}{38}{greyM}{navy}{잠금 대기}
\tfgb{-32}{38}{41}{navy}{white}{}
\tfgb{-32}{41}{73}{navyM}{white}{\code{Violations}}
\tfgb{-32}{73}{76}{navy}{white}{}
\node[note,anchor=north west] at (3,-34.3) {토큰 갱신 시점, 잠금 사용 중};
\node[note,anchor=north west] at (38,-34.3) {잠금 안에서 재확인: 토큰이 이미 새것, 두 번째 갱신 없음};
\node[note,anchor=west] at (77,-32) {반환};
% token lock
\tfgb{-42}{0}{36.5}{navy}{white}{R이 보유}
\node[note,anchor=west] at (38,-42) {비어 있음. G가 들어왔다 바로 나감};
% place lock
\tfgb{-51}{70}{145}{navy}{white}{검사부터 쓰기와 찾기까지 R이 보유}
% axis
\draw[greyD,line width=0.4pt] (0,-54.6) -- (145,-54.6);
\foreach \t in {0,35,70,105,140} {\draw[greyD,line width=0.4pt] (\t,-54.6) -- (\t,-55.6); \node[f,text=greyD,anchor=north] at (\t,-55.6) {\t};}
\node[f,text=greyD,anchor=north west] at (0,-58.8) {R과 G가 함께 시작한 뒤의 ms. 요청은 32\,ms로 그림(측정 31\textasciitilde 39\,ms). 진한 조각은 3\,ms 폭으로 그림};
\end{scope}
% ====================================================== panel (b): 0 to 10 s, 12 mm per s
\begin{scope}[yshift=-77mm]
\node[ptitle] at (-30,9) {(b) 처음 10\,s: 주문마다 피드 상태에서 찾아 확인한다. 쓰기 관문이 주문 2건의 간격을 벌린다};
\def\s{12}
\node[lane] at (-2,0)   {\textbf{피드 태스크}\\{\color{greyD}\code{hub.stream()}}};
\node[lane] at (-2,-12) {\textbf{코루틴 R}\\{\color{greyD}\code{rebalance()}}};
\node[lane] at (-2,-26) {\textbf{주문 잠금}};
% feed: sleeping, with request-reply slivers
\fill[navyL] (0,-2) rectangle (10*\s,2);
\foreach \t/\l in {1.88/{응답이 주문 1 추가},4.88/{응답이 주문 2 추가},7.88/{응답}} {
  \fill[navyM] (\t*\s-0.4,-2) rectangle (\t*\s+0.4,2);
  \node[f,anchor=south,text=navy] at (\t*\s,2.3) {\l};}
\foreach \t in {0.4,3.38,6.38} {\node[f,text=navy] at (\t*\s+5,0) {\code{sleep(3)}};}
\node[f,text=navy] at (8.88*\s+5,0) {\code{sleep(3)}};
% R lane
\fill[navy] (0,-14) rectangle (0.14*\s+0.6,-10);
\fill[navyL] (0.14*\s+0.6,-14) rectangle (2.14*\s,-10);
\fill[navy] (2.14*\s,-14) rectangle (2.25*\s+0.6,-10);
\fill[navyL] (2.25*\s+0.6,-14) rectangle (5.25*\s,-10);
\foreach \t in {0.64,1.14,1.64,2.75,3.25,3.75,4.25,4.75} {\draw[white,line width=0.7pt] (\t*\s,-14) -- (\t*\s,-10);}
\node[f,text=navy,anchor=west] at (5.25*\s+1.5,-12) {\code{rebalance()}가 주문 ID 2개를 반환};
% found markers and causal arrows
\foreach \t/\r in {2.14/1.88,5.25/4.88} {
  \fill[navy] (\t*\s,-9.6) -- ++(-1.1,1.8) -- ++(2.2,0) -- cycle;
  \draw[ev] (\r*\s,-2.2) -- (\t*\s,-7.6);}
\node[f,text=navy,anchor=west] at (2.14*\s+1.5,-6.6) {주문 1 찾음(2.0\,s)};
\node[f,text=navy,anchor=west] at (5.25*\s+1.5,-6.6) {주문 2 찾음(3.0\,s)};
% annotations below R
\draw[decorate,decoration={brace,amplitude=1.2mm,mirror},greyD,line width=0.5pt] (0.14*\s,-15) -- (1.14*\s,-15);
\node[note,anchor=north west,align=left] at (0,-16.8) {주문 1 보냄.\\관문 1\,s 닫힘};
\node[note,anchor=north west,align=left] at (2.14*\s-0.5,-15.4) {주문 2 보냄:\\관문 열림(2.1\,s)};
\node[note,anchor=north west,align=left] at (5.25*\s+1.5,-15.4) {흰 눈금: R이 최대 15\,s 동안 0.5\,s마다 \code{hub.state}를 보고,\\다음으로 REST \code{Orders}를 본다. 둘 다 없으면 R은\\\code{AmbiguousState}를 던지며 주문을 다시 보내지 않는다};
% place lock
\fill[navy] (0.07*\s,-28) rectangle (2.14*\s,-24);
\fill[navy] (2.14*\s,-28) rectangle (5.25*\s,-24);
\draw[white,line width=0.7pt] (2.14*\s,-28) -- (2.14*\s,-24);
\node[f,text=white] at ({(0.07+2.14)/2*\s},-26) {주문 1};
\node[f,text=white] at ({(2.14+5.25)/2*\s},-26) {주문 2: 검사, 쓰기, 찾기};
% axis
\draw[greyD,line width=0.4pt] (0,-31) -- (10*\s,-31);
\foreach \t in {0,...,10} {\draw[greyD,line width=0.4pt] (\t*\s,-31) -- (\t*\s,-32); \node[f,text=greyD,anchor=north] at (\t*\s,-32) {\t};}
\node[f,text=greyD,anchor=north] at (5*\s,-35.2) {R 시작 후 경과 시간(s)};
\end{scope}
% ====================================================== legend
\begin{scope}[yshift=-120mm]
\foreach \c/\x in {navy/0,navyM/34,greyM/64,navyL/96} {\fill[\c] (\x,0) rectangle ++(3.5,2.4);}
\node[f,anchor=west] at (4.3,1.2) {스레드에서 실행 중};
\node[f,anchor=west] at (38.3,1.2) {응답 대기};
\node[f,anchor=west] at (68.3,1.2) {잠금 대기};
\node[f,anchor=west] at (100.3,1.2) {잠자기 또는 상태 확인};
\fill[navy] (1.75,-2.6) -- ++(-1.1,1.8) -- ++(2.2,0) -- cycle; \node[f,anchor=west] at (4.3,-1.7) {주문 ID 찾음};
\draw[ev] (34,-1.7) -- (37.5,-1.7); \node[f,anchor=west] at (38.3,-1.7) {피드 응답이 \code{hub.state}를 바꿈};
\fill[navy] (96,-2.9) rectangle ++(3.5,2.4); \node[f,anchor=west] at (100.3,-1.7) {잠금 보유(잠금 줄)};
\end{scope}
\end{tikzpicture}}{잠금 줄에서 남색은 ‘보유’를 뜻한다. 측정값: Python 요청은 31\textasciitilde 39\ms{}가 걸리고(\ref{sec:transport}절), 주문은 \code{AddOrder} 후 1\textasciitilde 3\,s 만에 피드에 나타났다. 래퍼 설정: 3\,s마다 피드 요청, 0.5\,s마다 상태 확인, 쓰기 사이 1\,s, 하루 80건, 서버 시계로 720\,s 이후 첫 호출에서 갱신. 예시용: 시작 시각, 코루틴 순서, 피드 응답 시각은 측정 범위 안에서 골랐다. 간트 형식은 Silberschatz 외[7]를 따랐다. 패널 (b)에서 괄호 안 시간은 주문을 보낸 때부터 찾은 때까지, 그리고 마지막 쓰기가 끝난 때부터 다음 쓰기까지다.}{DART Quant Team; 코드는 \code{src/pytimefolio} (\code{auth}, \code{orders}, \code{hub})}


\lead{코드는 모듈 8개로 나뉜다. \ref{sec:layers}장의 5개 계층마다 하나씩, 피드 레코드를 위한 \code{state}, 이들을 묶는 \code{client}, 그리고 \code{errors}다(\figref{modules}).} 각 모듈은 표에서 자기보다 위에 있는 모듈만 가져오므로, 전송과 인증은 실시간 피드 없이, 주문 계산은 서버 없이 테스트할 수 있다. \code{errors}는 실패마다 예외를 하나씩 정의하며, 모두 \code{TimefolioError}의 하위 클래스다. 401은 \code{AuthError}, 400은 \code{RequestError}, 5xx나 끊긴 연결은 \code{ServerError}, 403이나 429는 \code{BlockedError}이고, 그 밖에 \code{StaleFeed}, \code{RebalanceError}, 그리고 목표 G5를 위한 \code{AmbiguousState}가 있다.

% Korean edition of figs/modules.tex.
\fig[wide]{modules}{모듈 8개: 프로토콜 계층마다 하나씩, 그리고 피드 상태, 클라이언트 클래스, 예외}[\code{pytimefolio}의 모듈 구성, import 순서]{%
{\tabfont\begin{tabularx}{\linewidth}{>{\RaggedRight\arraybackslash}p{18mm} >{\RaggedRight\arraybackslash}X >{\RaggedRight\arraybackslash}p{24mm}}
\thd{모듈} & \thd{역할} & \thd{근거}\\
\code{errors} & \code{TimefolioError}와 하위 클래스: \code{AuthError}(401), \code{RequestError}(400), \code{ServerError}(5xx나 끊긴 연결, 요청이 나갔는지 표시하는 플래그 포함), \code{BlockedError}(403이나 429), \code{StaleFeed}, \code{RebalanceError}(이미 넣은 ID 포함), \code{AmbiguousState} & \ref{sec:design}\\\rowsep
\code{transport} & 모든 REST 요청에 재사용하는 \code{httpx.AsyncClient(http2=True)} 하나. \code{warm_up()}. 응답 \code{Date} 헤더로 구한 서버 시계 오차. 상태 코드를 예외로 바꾸며, 어떤 요청에서든 처음 받은 403이나 429를 기록. GET은 연결이나 읽기/쓰기 실패 뒤 한 번 반복(타임아웃 뒤는 아님), POST는 절대 반복 안 함 & \ref{sec:transport}, \ref{sec:auth}\\\rowsep
\code{auth} & 시작할 때 저장된 쿠키 묶음으로 갱신, 실패할 때만 로그인. 갱신 나이를 넘긴 첫 호출이나 401에서 지연 갱신, 잠금 아래 한 번에 하나씩, 안에서 재확인. 쿠키 묶음 만료가 가까우면 로그인. 쿠키 묶음은 원자적으로, 소유자만 읽을 수 있게 기록. 비밀번호는 호출자에게서 받고 저장 안 함. 모든 호출에 Bearer 헤더 & \ref{sec:auth}\\\rowsep
\code{rest} & \code{currBizDay}와 KST 날짜 대조. 포트폴리오 규칙. \code{Session}, \code{Universe}, \code{Trades}, \code{Violations}, \code{Orders}(\code{d0}, \code{d1}), \code{ValidateProduct} & \ref{sec:queries}, \ref{sec:orders}\\\rowsep
\code{state} & 스냅샷 교체. \code{int(Id)}를 키로 한 델타 병합(\code{bs_updated}를 \code{bs}에 포함). 널 응답. \code{prodId} 색인. 가이드라인용으로 그 밖의 최상위 필드(\code{cmpl}, \code{tags}) 보관 & \ref{sec:feed}, \ref{sec:guidelines}\\\rowsep
\code{hub} & SignalR 협상, 토큰을 헤더에 실은 WebSocket, 핸드셰이크, 핑으로 연장되지 않는 20\,s 마감 하나를 가진 요청 하나, 핑에 퐁 응답. 마지막 응답 시각, 26\,s보다 오래되면 \code{StaleFeed}. 실패하면 지수 백오프와 full jitter로 재연결. 단 403이나 429면 피드 종료 & \ref{sec:feed}\\\rowsep
\code{orders} & 크기를 한 번 맞춘 주문 본문. water-filling, 섹터 검사, 진행 중 주문을 뺀 목표-주문 변환. 하루 한도(영업일마다 건수 파일), 쓰기 간격, 403이나 429 뒤 중단을 가진 주문 관문. 사전 검사(신선한 피드, 영업일, \code{ValidateProduct}, \code{weiLimit}). 쓰기와 찾기 내내 잠금 하나. 이미 돌려준 ID 건너뛰기, 새 레코드가 둘 이상 맞으면 \code{AmbiguousState}. 피드에서, 다음으로 REST \code{Orders}에서 확인하며 서버가 받았을 수 있는 주문은 재전송 안 함. 끝난 주문을 빼고 \code{CancelOrderBatch}로 정지. \code{Done}일 때 결과 판정 & \ref{sec:orders}, \ref{sec:oneshare}, \ref{sec:guidelines}\\\rowsep
\code{client} & 비동기 컨텍스트 매니저 \code{Timefolio}: 시작(세션, 포트폴리오, 피드 태스크, \code{Universe}, 오늘 주문 건수), 상태 보기, \code{buy}, \code{sell}, \code{stop}, \code{wait}, \code{nav()}, \code{share_weight()}, \code{plan()}, \code{rebalance()}, \code{guidelines()}, \code{orders_left()} & \ref{sec:design}, \ref{sec:live}\\
\end{tabularx}}}{각 모듈은 자기보다 위에 있는 모듈만 가져온다. 그 밖의 수치 매개변수(간격, 타임아웃, 한도, 갱신 나이)는 부록~\ref{app:decisions}에 있다. water-filling: 목표가 종목의 \code{weiLimit}를 넘으면 초과분을 다른 목표들에 비중에 비례해 나누므로 합계는 변하지 않는다.}{DART Quant Team; 코드는 \code{src/pytimefolio}}


\par\Needspace{6\baselineskip}
\claim{래퍼 운영: 실패마다 이름 붙은 예외를 던지고, 로그에는 개인 정보가 없다}

\lead{운영자, 즉 전략을 실행하는 사람이 개입해야 할 때를 알려주는 신호는 6개다. 예외 4개(\code{StaleFeed}, \code{BlockedError}, \code{AmbiguousState}, \code{RebalanceError}), 하루 건수, 가이드라인 플래그 \code{guidelines()["ok"]}다.} 아래 문단에서 차례로 다룬다. \code{StaleFeed}는 피드의 마지막 응답이 26\,s보다 오래됐을 때 상태 읽기, 주문 넣기, \code{wait()}이 던진다. 피드 태스크는 full-jitter 백오프로 계속 다시 연결하므로(\ref{sec:feed}절), 이 예외를 잡은 호출자는 피드가 다시 응답하면 재시도할 수 있다. 서버 핑은 계속 오는데 응답만 사라진 경우도 다른 타임아웃과 똑같이 20\,s 뒤에 끝난다. 핑은 요청의 마감 시간을 늘리지 않기 때문이다.

\lead{어떤 요청에서든 403이나 429를 받으면 \code{BlockedError}를 던지고, 프로세스가 끝날 때까지 모든 쓰기를 멈춘다.} 주최측은 트래픽 폭주를 이유로 IP 주소를 차단할 수 있고, 래퍼가 볼 수 있는 차단 신호는 403이나 429뿐이다. 전송 계층은 읽기, 쓰기, 피드 협상 등 어떤 요청에서든 처음 받은 403이나 429를 기록하고, 주문 관문은 이후 모든 쓰기를 거부한다. 협상에서 차단되면 피드도 다시 연결하지 않는다. 피드 태스크가 끝나고, 이후 상태 읽기와 주문은 차단을 명시한 예외를 던진다. 다른 요청에서 차단된 경우에는 피드가 계속 돌아가므로 호출자가 상태를 기록할 수 있다. 다만 IP 차단이면 읽기도 거부될 수 있다. 주문은 다시 시작해야만 재개되며, 그 전에 사람이 주최측에 먼저 문의할 기회가 생긴다.

\lead{래퍼가 확인할 수 없는 주문은 \code{AmbiguousState}를 던지고 다시 보내지 않는다(\ref{sec:orders}절). 실패한 \code{rebalance()}는 \code{RebalanceError}를 던진다.} \code{rebalance()}는 첫 번째 실패한 주문에서 멈추며, 예외의 \code{placed} 속성에 이미 보낸 주문 ID가 들어 있어 호출자가 그 주문을 기다리거나 정지할 수 있다. 설계는 최대 한 번(at-most-once)이다. 모호한 응답은 모두 운영자가 확인할 일이 되며, 이는 주문을 두 번 넣지 않기 위해 치르는 비용이다.

\lead{하루 건수는 80 이하로 유지되며, 시작할 때와 \code{rebalance()} 때마다 서버의 건수로 올린다.} \code{AddOrder}마다 1을 더하고, 정지한 주문마다 1을 더한다. \code{CancelOrderBatch}는 정지하는 ID를 모두 세기 때문이다. 서버에 닿지 않은 요청은 더하지 않는다. 건수는 \code{~/.config/}\allowbreak\code{pytimefolio} 아래 영업일마다 파일 하나에 파일 잠금을 잡고 원자적으로 기록한다. 시작할 때와 \code{rebalance()} 때마다 로컬 건수와 서버에 기록된 그날의 새 주문·정지 건수 중 큰 값으로 올리므로, 브라우저에서 넣은 주문도 센다. 11월 24일 10:34 기준 카운터는 80건 중 16건을 썼다(새 주문 9건, 정지 7건). 80을 넘길 요청은 보내기 전에 예외를 던진다. 한 프로세스 안에서는 한도가 정확하다. 건수 파일을 공유하는 여러 프로세스 사이에서는 쓰기 전 검사가 파일 잠금을 잡지 않으므로 프로세스 수만큼의 오차 안에서만 지켜진다.

\lead{\code{guidelines()["ok"]}는 서버가 센 현재 위반이 0건이고 이번 주 회전율이 5\% 이상일 때만 True다(\ref{sec:guidelines}절).} 따라서 호출자는 리밸런싱 전후에 플래그 하나만 확인하면 된다.

\lead{로그는 래퍼가 한 일을 기록하며, 자격 증명과 개인 정보 필드는 모두 뺀다.} 래퍼는 Python \code{logging} 모듈로 로거 \code{pytimefolio}와 하위 로거 \code{auth}, \code{hub}, \code{orders}에 기록하며, 로그를 어디로 보낼지는 호출자가 정한다. README의 예제는 로그를 표준 오류로 보내고 \code{httpx}는 WARNING으로 올린다. \code{httpx}는 INFO 수준에서 모든 URL을 기록하기 때문이다. 래퍼는 시작 줄(영업일, 80건 중 사용 건수), 로그인마다 쿠키 만료 시각, 갱신마다 영업일, 피드 연결과 실패마다 백오프 대기 시간, 넣은 주문마다(주문 ID, 종목코드, \code{wei}, \code{state}, \code{qty2ord}), 그리고 차단 신호를 기록한다. 개인 정보(연락처)가 들어 있는 로그인 응답, 포지션 레코드의 이름(\code{userNm}, \code{userNick}), 액세스 토큰, 쿠키, 비밀번호는 절대 기록하지 않는다. 토큰은 헤더로 전달되므로 \code{httpx}가 기록하는 URL에도 들어가지 않는다.

\lead{쿠키 묶음은 래퍼가 디스크에 쓰는 유일한 자격 증명이므로, 래퍼는 \code{session.json}을 7일짜리 자격 증명으로 다룬다(\figref{security}).} 비밀번호는 생성자로 넘기며(예: 환경 변수에서) 저장하지 않는다. 쿠키 묶음은 저장소 밖의 \code{~/.config/}\allowbreak\code{pytimefolio/}\allowbreak\code{session.json}에 권한 0600으로, 로그인과 갱신 때마다 원자적으로 기록한다. 이 파일은 브라우저에서 \code{Logout}을 해도 최대 7일 동안 계속 갱신된다(\ref{sec:auth}절).

% Korean edition of figs/security.tex.
\fig[wide]{security}{쿠키 묶음은 래퍼가 저장하는 유일한 비밀이다}[\code{pytimefolio}를 거치는 비밀과 개인 정보, 그리고 각각을 다루는 방식]{%
{\tabfont\begin{tabularx}{\linewidth}{>{\RaggedRight\arraybackslash}p{30mm} >{\RaggedRight\arraybackslash}p{62mm} >{\RaggedRight\arraybackslash}X}
\thd{항목} & \thd{허용하는 것 또는 담고 있는 것} & \thd{래퍼가 다루는 방식}\\
\hl 쿠키 묶음 & 로그인 후 7일간 액세스 토큰을 갱신. \code{Logout} 뒤에도 계속 갱신됨. 비밀번호 변경 뒤 동작은 미시험 & 저장소 밖 \code{~/.config/pytimefolio/session.json}에 권한 0600으로 저장. 로그인과 갱신 때마다 원자적으로 기록\\\rowsep
비밀번호 & 로그인하고 새 7일짜리 쿠키 묶음을 만듦 & 래퍼가 저장하지 않음. 호출자가 생성자에 넘김(예: 환경 변수에서)\\\rowsep
액세스 토큰 & 900\,s 동안 모든 REST·피드 요청을 인증 & 메모리에만 두고 \code{Authorization} 헤더로만 보냄. URL과 로그에 남기지 않음\\\rowsep
로그인 응답 & 개인 정보 필드(연락처) & 기록하지 않음\\\rowsep
포지션 레코드 & 계정 소유자의 이름과 닉네임(\code{userNm}, \code{userNick}) & 기록하지 않음\\\rowsep
저장소 & GitHub에 공개 & \code{.env}, \code{data/}, \code{logs/}는 \code{.gitignore}에 있음\\
\end{tabularx}}}{음영: 디스크에 쓰는 유일한 비밀. 토큰과 쿠키의 유효기간은 \ref{sec:auth}절에서 가져왔다. 비밀번호 변경 테스트는 \ref{sec:validity}절의 미해결 질문에 있다.}{DART Quant Team; 코드는 \code{src/pytimefolio}}

